\section{约定}\label{sec:preliminaries}

\paragraph{带权博弈。}
中间实例可以有自环、重复的约束出现、以及和为 1 的正有理权重。其值
是某个染色所满足的最大总权重。即使约束重合，各出现仍保留独立的
ID。所有有理数都用二进制分子与分母编码。

二部\emph{投影博弈}在两侧每个顶点上有一个问题、有限的答案集，以及
每条约束上从左答案到右答案的映射。策略为每个问题选一个答案。依赖
问题的答案集可以与同规模的固定字母表等同；对奇偶方程的四个满足
三元组我们就这么用。每个证明者只用自己的问题。与问题独立的私有
或共享随机性给出确定性策略的混合，不会增大值。

并行幂 \(G^{\otimes t}\) 采样 \(t\) 个独立的约束出现，发送有序的
问题元组，当每个坐标都接受时接受。答案可以依赖证明者的整个元组。
重复的问题 ID 保留独立的坐标位置与独立的局部答案副本。这就是重复
定理所用的普通经典乘积博弈。

\paragraph{二元线性代数与傅里叶分析。}
除特别指明更大域外，所有向量空间与映射都是 \(\F_2\) 上的有限维。
记 \(\Hom(E,F)\) 为线性映射空间，\(E^*=\Hom(E,\F_2)\) 为对偶。
子空间 \(S\le E^*\) 在子空间 \(K\le E\) 上的限制为其像
\(S|_K\le K^*\)；其限制秩为 \(\dim(S|_K)\)。

对二元空间 \(E\) 上的实函数 \(f\)，使用归一化傅里叶系数
\[
 \widehat f(g)=\E_{x\in E}f(x)(-1)^{g(x)},\qquad
 f(x)=\sum_{g\in E^*}\widehat f(g)(-1)^{g(x)},\qquad
 \sum_g|\widehat f(g)|^2=\E_x|f(x)|^2.
\]
这些恒等式源于：非零线性泛函取零与取一的次数相同，故其特征均值
为零。因此这些特征构成正交规范基。在矩阵空间上，这是与其对偶矩阵
空间的逐项配对；定义傅里叶解码器时会指明张量记号。

\paragraph{概率与有限构造。}
从向量空间均匀采样包括零。除非声明依赖，样本相互独立。对有限概率
律，
\[
 \TV(P,Q)=\frac12\sum_x|P(x)-Q(x)|.
\]
它控制任何 \([0,1]\) 值函数的期望之差。我们只在正概率事件上作
条件化。

随机实验描述有限的带权约束列表。归约确定性地枚举这些列表。每个
辅助维数与元组长度都在输入公式变化之前固定；最终证明同时检查该
参数顺序与二进制编码代价。傅里叶见证与解码器只用于界定博弈值，
不必高效。
